% s-305-p305-m002; original PDF 305, 306
\Aphorism{s-305-p305-m002}{359}
Es giebt, soweit überhaupt \Foreign{factum infectum fieri quit,} eine einfach menschliche Art, Verschuldungen wieder gut zu machen, eine Art, die sich alle voreilige Verallgemeinerung, alle tragischen Gebärden der Schuld, alle Entfaltung grosser moralischer Decorationsstücke verbietet und es nicht verschmäht, den thatsächlichen Bestand der Schuld, die oft verwunderlich geringe \Foreign{injuria} im Sinne des Geschädigten, zu tilgen. Bisweilen heilt man eine Wunde, die man schlug, mit freundlichen Worten, sogar mit Geld oder noch harmloseren Mitteln — das ist nun freilich nichts für die Komödianten des Schuld- und Sühneverhältnisses, für die Liebhaber schwerfälligen Opfer- und Vergeltungspompes, für Jene, die irgend einen Theil von sich, sei es auch nur eine Verschuldung, ins Ungeheuerliche, Wolkenhafte, Ewige, Unwiderrufliche auffärben müssen, um sich interessant zu erscheinen. Von solchen gewohnheitsmässigen Hyperbolikern, die im Purpurmantel der Schuld nicht ohne Coquetterie einherwandeln, wird mit ärgerlicher Einseitigkeit nur die \Foreign{active injuria}, die Schuld mit Beziehung auf den Schuldigen allein, empfunden, während sie den Geschädigten als blosses zufälliges Object bei Seite schieben; \Sperr{ihr} Seelenzustand, \Sperr{ihre} Rechtfertigung, die transcendente Vertiefung und Verteufelung \Sperr{ihrer} Wenigkeit — damit erhitzt sich ihre Phantasie bis zum piquanten Unsinn. Es kommt vor, dass die künstlich aufgeblasene Schuld zu kosmischer Grösse anschwillt; die verletzte „sittliche Weltordnung“, ein welthistorischer Frevel, Kain, Judas, Herodias, lauter verzerrte Ungeheuer aus einer moral-mythologischen \Foreign{laterna magica} erscheinen an der Wand. Einige Moderne übergaukeln mit solchem Fratzenwesen noch die christlichen Dogmatiker; Schopenhauers Schuldbegriff (in Dostojewskij’scher Formel „Schuld an Allem und für Alle“) lässt die Erbsünde tief unter sich, und wenn Parsifal im dritten Act aufschreit „ich bin’s, der all dies Elend schuf!“ — so erstaunt man, milde geredet, über die Folgen eines einzigen Kusses. — Wenn die Infection mit moralischen Nebenabsichten uns das Theater verdorben hat, so hat doch auch umgekehrt das Theatralische die Moral verdorben. Jede dieser beiden Volksthümlichkeiten hat die andere um ihre specielle Untugend bereichert: das Theater die Moral um seine gespreizten Attitüden, die Moral das Theater um ihren beschränkten Horizont.
\par

